\section{Motivación}
La segmentación automática de vasos retinianos es un componente crítico en sistemas de apoyo al diagnóstico de enfermedades oftalmológicas y cardiovasculares. Un delineado preciso de la red vascular permite extraer biomarcadores estructurales y hemodinámicos útiles para tamizaje clínico y monitoreo longitudinal.

Sin embargo, el desempeño de los modelos puede degradarse cuando se evalúan en distribuciones diferentes a las de entrenamiento. Este fenómeno es especialmente relevante en retina, donde cambian condiciones de iluminación, resolución, contraste, sensor y criterios de anotación entre conjuntos de datos.

En este contexto, el objetivo del trabajo es doble: (1) obtener un modelo robusto en distribución sobre DRIVE y (2) cuantificar y analizar su capacidad de generalización sobre CHASE\_DB1. A partir de ello, se propone una estrategia de adaptación de dominio por preprocesamiento para reducir la brecha de rendimiento entre conjuntos.
